\documentclass[11pt]{article}

\usepackage[margin=1in]{geometry}
\usepackage{amsmath,amssymb,amsthm,mathtools}
\usepackage{stmaryrd}
\usepackage{xcolor}
\usepackage{hyperref}
\usepackage{tikz}
\usetikzlibrary{positioning,arrows.meta,calc,fit,backgrounds}
\usepackage{enumitem}

\hypersetup{
  colorlinks=true,
  linkcolor=blue!50!black,
  citecolor=blue!50!black,
  urlcolor=blue!50!black
}

% --- Theorem environments ---
\theoremstyle{plain}
\newtheorem{theorem}{Theorem}[section]
\newtheorem{lemma}[theorem]{Lemma}
\newtheorem{proposition}[theorem]{Proposition}
\newtheorem{corollary}[theorem]{Corollary}

\theoremstyle{definition}
\newtheorem{definition}[theorem]{Definition}
\newtheorem{example}[theorem]{Example}
\newtheorem{construction}[theorem]{Construction}

\theoremstyle{remark}
\newtheorem{remark}[theorem]{Remark}

% --- Notation ---
\newcommand{\sem}[1]{\llbracket\,#1\,\rrbracket}
\newcommand{\reflect}[1]{\ulcorner #1 \urcorner}
\newcommand{\rew}{\leadsto}
\newcommand{\Rew}{\Longrightarrow}
\newcommand{\Names}{\mathcal{N}}
\newcommand{\Procs}{\mathcal{P}}
\newcommand{\Terms}{\mathcal{T}}
\newcommand{\Pos}{\mathbb{P}}
\newcommand{\States}{\mathcal{S}}
\newcommand{\Lhs}{\mathcal{L}}
\newcommand{\Patterns}{\mathcal{L}}
\newcommand{\hd}{\mathsf{hd}}
\newcommand{\arity}[1]{\##1}
\newcommand{\dom}{\mathcal{D}}
\newcommand{\edges}{\mathcal{E}}
\newcommand{\rdep}{R_{\mathsf{dep}}}
\newcommand{\rop}{R_{\mathsf{op}}}
\newcommand{\gcp}{\mathsf{gcp}}
\newcommand{\reduce}{\mathsf{reduce}}
\newcommand{\deriv}{\mathsf{deriv}}
\newcommand{\lift}{\mathsf{lift}}
\newcommand{\surface}{\mathsf{surf}}
\newcommand{\holes}{\mathsf{holes}}
\newcommand{\true}{\mathsf{true}}
\newcommand{\false}{\mathsf{false}}
\newcommand{\Match}{\mathit{match}}
\newcommand{\nil}{\mathbf{0}}

% Short rho-syntax
\newcommand{\Send}[2]{#1 \,!\, (#2)}
\newcommand{\Recv}[3]{\mathsf{for}\,(#1 \Leftarrow #2)\,\{\,#3\,\}}
\newcommand{\Let}[3]{\mathsf{let}\;#1 = #2\;\mathsf{in}\;#3}

\title{Optimal Channel Naming for Compositional Rewrite Translations\\
       via Set Automaton Partial Evaluation}
\author{L.G.~Meredith\\\small F1R3FLY.io}
\date{\today}

\begin{document}
\maketitle

\begin{abstract}
We present a translation scheme that compiles graph-structured lambda-theory
(GSLT) rewrites into the rho calculus, and we identify the channel name for
each rewrite context as the canonical reflection of a state in a Bouwman--Erkens
set automaton. Treating the polymorphic compilation step
$\Let{tc}{\sem{K}}{\cdots}$ as a partial evaluation of the set automaton
against the context $K$ yields a closed-form formula
$tc(K) = \reflect{\delta^{*}(s_{0}, \surface(K))}$. We then transport the
three principal theorems of the Bouwman--Erkens construction --- symbol-once
inspection, prune-preserves-work, and the coarsest-sound-partition theorem
for the outermost-preserving dependency relation --- into the channel-based
setting. The resulting scheme is optimal in three precise senses:
each surface symbol of every redex context is consumed by exactly one
\textsf{for}-receive, inner reductions never invalidate outer channels, and
the channel quotient induced by $tc(\cdot)$ is the coarsest equivalence on
contexts that preserves outermost firing. We discuss the non-linear
extension via rho-calculus consistency receives, and outline the changes
required to support dynamic rule sets.
\end{abstract}

\tableofcontents

\section{Introduction}\label{sec:intro}

Compiling a rewrite system into a process calculus poses a tension that does
not arise when implementing a rewrite engine in an imperative or functional
host language. In the latter setting, a rewrite procedure can freely re-traverse
the term being rewritten, and an efficient pattern matcher --- such as the
set automaton of Bouwman and Erkens~\cite{BouwmanErkens2022,ErkensGroote2021}
--- ensures that each function symbol is inspected only once per redex search.
In the former setting, every traversal step is a communication, every
intermediate matching state must be reified as a channel or process, and
every redex firing must be expressible as a send-receive pair.

The translation scheme considered here, due to Meredith, compiles a single
rewrite rule $L \rew R$ as
\begin{equation}\label{eq:single-rule}
\sem{L \rew R}(t\!\ell)
\;\;=\;\;
\Recv{\sem{L}}{t\!\ell}{\,\Send{t\!\ell}{\sem{R}}\,}
\end{equation}
and a contextual rewrite as
\begin{equation}\label{eq:context-rule}
\begin{array}{@{}l@{}}
\sem{L_{1} \rew R_{1}, \ldots, L_{n} \rew R_{n}
     \;\Rightarrow\; K(L_{1}, \ldots, L_{n}) \rew K'(R_{1}, \ldots, R_{n})} \;=\; \\[2pt]
\quad \mathsf{let}\;tc = \sem{K}\;\mathsf{in} \\
\quad\quad \mathsf{for}\,((\sem{L_{1}}, \ldots, \sem{L_{n}}) \Leftarrow tc)\,\{ \\
\quad\quad\quad \Send{tc}{\sem{K'}(\sem{R_{1}}, \ldots, \sem{R_{n}})} \\
\quad\quad \}
\end{array}
\end{equation}
The polymorphic interpretation of $\Let{tc}{\sem{K}}{\cdots}$ is the locus
of optimisation. Since $K$ is a context (a term with $n$ holes), the channel
$tc$ is a function of $K$'s shape; the question is which function makes the
overall translation optimal.

We argue that the right answer is given by partial evaluation of a
set automaton built from the left-hand sides $\{L_{1}, \ldots, L_{n}\}$.
Concretely:
\[
  tc(K) \;=\; \reflect{\delta^{*}(s_{0}, \surface(K))},
\]
where $(\States, s_{0}, L, \delta, \eta)$ is the set automaton for the
pattern set, $\surface(K)$ is the function-symbol skeleton of $K$ in the
inspection order prescribed by the state-labelling function $L$, and
$\reflect{\cdot}$ is the rho-calculus reflection of an automaton state into
a channel name. Two contexts $K$ and $K'$ that lead the automaton to the
same state are then assigned the same channel; two contexts that the
automaton distinguishes are assigned distinct channels. We show that the
partition on contexts induced by $tc(\cdot)$ is the coarsest partition
sound for outermost firing.

\paragraph{Contributions.}
\begin{itemize}[leftmargin=*]
  \item A formal definition of the channel-naming function $tc(K)$ as a
        partial evaluation of a set automaton (\S\ref{sec:construction}).
  \item Three optimality theorems transported from~\cite{BouwmanErkens2022}:
        symbol-once inspection (Theorem~\ref{thm:symbol-once}),
        prune-preserves-work (Theorem~\ref{thm:prune-preserves}), and
        coarsest sound partition (Theorem~\ref{thm:coarsest}).
  \item A treatment of the non-linear case
        (\S\ref{sec:nonlinear}) via consistency receives, paralleling
        the Bouwman--Erkens \emph{ambiguous/disabled/enabled} bookkeeping.
  \item A discussion of dynamic rule sets and other engineering caveats
        (\S\ref{sec:caveats}).
\end{itemize}

\paragraph{Outline.}
Section~\ref{sec:prelim} fixes notation for the rho calculus, GSLT rewrites,
and set automata. Section~\ref{sec:translation} reviews the translation
scheme and isolates the polymorphic step.
Section~\ref{sec:construction} gives the channel-computation algorithm.
Section~\ref{sec:optimality} proves the optimality theorems.
Section~\ref{sec:nonlinear} treats non-linearity.
Section~\ref{sec:caveats} discusses caveats and future work.

\section{Preliminaries}\label{sec:prelim}

\subsection{Terms, contexts, and rewrites}

Fix a finite ranked alphabet $\mathbb{F}$ with arity map
$\arity{\cdot} : \mathbb{F} \to \mathbb{N}$ and a countable set of
variables $V$, with $\mathbb{F} \cap V = \emptyset$. The set
$\Terms(\mathbb{F}, V)$ of terms is generated by $V \subseteq
\Terms(\mathbb{F}, V)$ together with the formation rule that
$f(t_{1}, \ldots, t_{n}) \in \Terms(\mathbb{F}, V)$ whenever
$f \in \mathbb{F}$ has $\arity{f} = n$ and $t_{1}, \ldots, t_{n} \in
\Terms(\mathbb{F}, V)$. We write $\mathit{vars}(t)$ for the set of
variables occurring in $t$ and call $t$ \emph{ground} when
$\mathit{vars}(t) = \emptyset$.

A \emph{position} is a finite list of positive integers; we write
$p, q \in \Pos$, $\epsilon$ for the empty list (the \emph{root
position}), and $p \cdot q$ for concatenation, satisfying
$\epsilon \cdot p = p \cdot \epsilon = p$. The \emph{prefix order}
$p \le q$ holds iff $p \cdot r = q$ for some $r$. The \emph{domain} of a
term is defined by $\dom(x) = \{\epsilon\}$ for $x \in V$ and
$\dom(f(t_{1}, \ldots, t_{n})) = \{\epsilon\} \cup
\{\, i \cdot p \mid 1 \le i \le n,\; p \in \dom(t_{i}) \,\}$. For
$p \in \dom(t)$, the subterm at $p$ is $t|_{p}$, with
$t|_{\epsilon} = t$ and $f(t_{1}, \ldots, t_{n})|_{i \cdot p'} =
t_{i}|_{p'}$; the replacement $t[u]_{p}$ is the unique term obtained
from $t$ by replacing $t|_{p}$ with $u$. The \emph{edge set}
$\edges(t) \subseteq \dom(t)$ is the set of positions of variable
occurrences, and $\dom_{\setminus\edges}(t) = \dom(t) \setminus
\edges(t)$ is the set of function-symbol positions. For a non-variable
term $t = f(t_{1}, \ldots, t_{n})$ the \emph{head symbol} is
$\hd(t) = f$.

A \emph{pattern} is a non-variable term; a pattern is \emph{linear} if
no variable occurs in it more than once. A \emph{substitution} is a
mapping $\sigma : V \to \Terms(\mathbb{F}, V)$; we write $t^{\sigma}$
for its homomorphic extension applied to $t$.

A \emph{context} of arity $n$ is a term-with-holes
$K(\Box_{1}, \ldots, \Box_{n})$ in which each hole $\Box_{i}$ occurs
exactly once. We write $H_{K} \subseteq \dom(K)$ for the set of hole
positions and $K[t_{1}, \ldots, t_{n}]$ for the term obtained by
simultaneous substitution of $t_{i}$ for $\Box_{i}$.

A \emph{rewrite rule} $L \rew R$ is a pattern $L$ paired with a term $R$
such that $\mathit{vars}(R) \subseteq \mathit{vars}(L)$. A
\emph{contextual rewrite} is a sequent
\[
  L_{1} \rew R_{1}, \ldots, L_{n} \rew R_{n}
  \;\Rightarrow\;
  K(L_{1}, \ldots, L_{n}) \rew K'(R_{1}, \ldots, R_{n}),
\]
asserting that, in any context where the inner rewrites are
applicable, the outer pattern reduces to the outer right-hand side.

\paragraph{Binders.}
A function symbol $f \in \mathbb{F}$ of arity $n$ may be designated a
\emph{binder} at one or more argument positions
$\beta_{f} \subseteq \{1, \ldots, n\}$. Under the GSLT semantics
of~\S\ref{ssec:gslt}, a binding argument at position $i \in \beta_{f}$
has type $P \Rightarrow P$ via the internal hom of $\mathcal{C}$
rather than $P$; the canonical reading is HOAS in the host calculus.
For position-based reasoning we \emph{open} such a binder by
introducing a fresh free name: in $f(t_{1}, \ldots, t_{n})$ with
$i \in \beta_{f}$ the subterm at $p \cdot i$ is the opened body
\[
  t|_{p \cdot i} \;=\; t_{i}(z_{p \cdot i}),
\]
where $z_{p \cdot i}$ is a free name canonically associated with the
position $p \cdot i$. The $z_{p \cdot i}$ are pairwise distinct and
distinct from any variable in $V$ occurring in the surrounding term;
the term is determined only up to $\alpha$-renaming of those names.
The fresh free names are \emph{leaves}: $\dom(z_{p}) = \{\epsilon\}$
and $\hd(z_{p})$ is undefined, exactly as for variables in $V$.
The formulae for $\dom$, $\edges$, $t|_{p}$, and $\hd$ above extend
by the same recursion through opened bodies, and a presentation may
omit any first-order $\mathit{var}$ constructor: object-language
variables \emph{are} the host calculus's free names introduced at
binder openings.

\subsection{Graph-structured lambda theories}\label{ssec:gslt}

We adopt the semantic definition of a graph-structured lambda theory
and treat term-rewriting data as a finite \emph{presentation} of one.

\begin{definition}[Graph-structured lambda theory]\label{def:gslt}
A \emph{graph-structured lambda theory} (GSLT) is a tuple
$\mathcal{G} = (\mathcal{C}, P, R)$ where
\begin{itemize}[leftmargin=*]
  \item $\mathcal{C}$ is a category with finite limits;
  \item $\mathcal{C}$ is closed with respect to the cartesian product
        (i.e.\ for every object $A$ the functor $({-}) \times A$ has
        a right adjoint $A \Rightarrow ({-})$, giving internal homs);
  \item $P \in \mathcal{C}$ is a distinguished object, called the
        \emph{object of processes};
  \item $R \hookrightarrow P \times P$ is a distinguished subobject,
        called the \emph{reduction relation}.
\end{itemize}
A \emph{syntactic morphism} of GSLTs $\mathcal{G} \to \mathcal{G}'$
is a functor $F : \mathcal{C} \to \mathcal{C}'$ that strictly preserves
finite limits and internal homs, satisfies $F(P) = P'$ on the nose,
and restricts on $R$ to a map into $R'$ over $P' \times P'$.
\end{definition}

The reduction relation $R$ is given as an abstract subobject; in
practice the GSLTs of interest are obtained as the free GSLT generated
by a finite collection of generators and relations. Term rewriting
provides one such presentation.

\begin{definition}[Finite presentation]\label{def:presentation}
A \emph{finite presentation} of a GSLT consists of a finite ranked
alphabet $\mathbb{F}$ (and an apparatus for binders, if needed)
together with a finite set
\[
  \mathcal{R} \;=\; \mathcal{R}_{\mathrm{rule}}
                    \,\cup\, \mathcal{R}_{\mathrm{ctx}}
\]
of rewrite rules and contextual rewrites in the sense above.
\end{definition}

\begin{construction}[GSLT generated by a presentation]\label{constr:gslt}
Let $\mathcal{R}$ be a presentation over $\mathbb{F}$. The
\emph{generated GSLT}
$\mathcal{G}(\mathcal{R}) = (\mathcal{C}_{\mathbb{F}}, P, R(\mathcal{R}))$
is built as follows.
\begin{itemize}[leftmargin=*]
  \item $\mathcal{C}_{\mathbb{F}}$ is the cartesian closed category
        freely generated by the signature: objects are simple types
        built from a base sort $\iota$ under $\times$ and $\Rightarrow$,
        and morphisms are $\lambda$-terms modulo $\beta\eta$ over the
        function symbols of $\mathbb{F}$ (equivalently, the syntactic
        category of the simply-typed $\lambda$-theory whose constants
        are $\mathbb{F}$). This category has all finite limits and is
        cartesian closed by construction.
  \item $P$ is the base sort $\iota$; ground terms of $\Terms(\mathbb{F})$
        are the global elements $1 \to P$.
  \item $R(\mathcal{R}) \hookrightarrow P \times P$ is the
        \emph{inductively generated} subobject specified in
        Definition~\ref{def:reduction-relation} below.
\end{itemize}
\end{construction}

\begin{definition}[Inductively generated reduction relation]\label{def:reduction-relation}
$R(\mathcal{R}) \hookrightarrow P \times P$ is the smallest subobject
closed under the following clauses, read pointwise on global elements
(equivalently, on ground terms over $\mathbb{F}$):
\begin{description}[leftmargin=2em,style=nextline]
  \item[(Rule)] For every rewrite rule $L \rew R$ in
        $\mathcal{R}_{\mathrm{rule}}$ and every ground substitution
        $\sigma$, $(L^{\sigma}, R^{\sigma}) \in R(\mathcal{R})$.
  \item[(Context)] For every contextual rewrite
        $L_{1} \rew R_{1}, \ldots, L_{n} \rew R_{n} \Rightarrow
         K(L_{1}, \ldots, L_{n}) \rew K'(R_{1}, \ldots, R_{n})$
        in $\mathcal{R}_{\mathrm{ctx}}$ and every ground substitution
        $\sigma$: if
        $(L_{i}^{\sigma}, R_{i}^{\sigma}) \in R(\mathcal{R})$ for each
        $1 \le i \le n$, then
        $\bigl(K^{\sigma}(L_{1}^{\sigma}, \ldots, L_{n}^{\sigma}),\,
              K'^{\sigma}(R_{1}^{\sigma}, \ldots, R_{n}^{\sigma})\bigr)
         \in R(\mathcal{R})$.
\end{description}
Equivalently, $R(\mathcal{R})$ is the least fixed point of the
monotone operator on subobjects of $P \times P$ that adds the
$\sigma$-instances of the rule clauses and the $\sigma$-instances of
the context clauses whose premises already lie in the argument.
\end{definition}

\begin{remark}\label{rmk:presentation}
The GSLT proper is the generated $\mathcal{G}(\mathcal{R})$ of
Construction~\ref{constr:gslt}. By a mild abuse we write $\mathcal{R}$
for both the presentation and the GSLT it generates wherever no
confusion can arise; statements about ``redexes,'' ``firing,'' and
``the reduction relation'' should be read as statements about
$R(\mathcal{R})$.
\end{remark}

The set of \emph{left-hand sides} of a presentation $\mathcal{R}$ is
\[
  \Lhs \;=\; \{L \mid L \rew R \in \mathcal{R}_{\mathrm{rule}}\} \cup
  \{L_{i} \mid L_{1} \rew R_{1}, \ldots, L_{n} \rew R_{n} \Rightarrow \cdots \in \mathcal{R}_{\mathrm{ctx}},\; 1 \le i \le n\}.
\]
The set automaton of \S\ref{ssec:set-automata} is built from $\Lhs$
alone; it is the syntactic data of the presentation, and is the same
for any two presentations that share their left-hand sides.

\subsection{Rho calculus, briefly}

We use a minimal fragment of the rho calculus~\cite{MeredithRadestock2005}.
Processes $P$ and names $x$ are mutually defined by
\[
  P \;::=\; \nil \;\mid\; x!(P) \;\mid\; \mathsf{for}\,(y \Leftarrow x)\,\{P\}
       \;\mid\; P \mid P \;\mid\; *x,
  \qquad
  x \;::=\; @P,
\]
where $\nil$ is the inert process, $\Send{x}{P}$ is an asynchronous send
of $P$ on channel $x$, $\Recv{y}{x}{P}$ is a persistent receive on
channel $x$ that binds the incoming name to $y$ in $P$, $P \mid Q$ is
parallel composition, $@P$ is the \emph{reflection} of process $P$ to a
name, and $*x$ is the \emph{dereference} of name $x$ to a process.

The rho calculus has a structural congruence $\equiv$ that includes
$\alpha$-equivalence, associativity-commutativity of $\mid$ with $\nil$
as unit, and the reflection law $*@P \equiv P$. Names are quotients of
processes under $\equiv$: two names $x = @P$ and $x' = @P'$ are
\emph{name-equal}, written $x =_{\Names} x'$, iff $P \equiv P'$. The
principal reduction rule (COMM) is
\[
  \Recv{y}{x}{P} \;\mid\; \Send{x'}{Q}
  \;\rew\;
  P[@Q / y]
  \qquad
  \text{whenever } x =_{\Names} x',
\]
with the persistent variant leaving a fresh copy of the receive in place.

We use $\Let{x}{P}{Q}$ as syntactic sugar for $Q[@P/x]$ when $P$ is a
name-valued expression. Tuple-pattern receives
$\Recv{(y_{1}, \ldots, y_{n})}{x}{P}$ are the standard pattern-match on
incoming names, decomposable into a finite number of nested single-name
receives guarded by name equality (an explicit elaboration is unnecessary
for what follows).

The reflection bracket $\reflect{\cdot}$, used in the body of the paper, is
a section of $@\cdot$ chosen so that distinguishable mathematical objects
are mapped to distinct names: $\reflect{a} =_{\Names} \reflect{b}$ iff
$a$ and $b$ are equal as mathematical objects. For automaton states
$s \in \States$, $\reflect{s}$ denotes a canonical name unique to $s$;
we extend this convention to configuration trees in
\S\ref{sec:construction}.

\subsection{Set automata}\label{ssec:set-automata}

We summarise the construction of~\cite{BouwmanErkens2022,ErkensGroote2021}
in enough detail to make the rest of the paper self-contained. Readers
already familiar with their work may skim this subsection.

\paragraph{The tuple.}
A \emph{set automaton} for a finite pattern set $\Lhs$ over
$\mathbb{F}$ is a tuple
\[
  M \;=\; (\States,\, s_{0},\, L,\, \delta,\, \eta),
\]
where $\States$ is a finite set of \emph{states}, $s_{0} \in \States$
is the \emph{initial state}, $L : \States \to \Pos$ is the
\emph{state-labelling} function (it selects the next position to
inspect from a configuration in state $s$),
$\delta : \States \times \mathbb{F} \to \mathcal{P}(\States \times \Pos)$
is the \emph{transition} function, and $\eta : \States \times \mathbb{F}
\to \mathcal{P}(\Lhs \times \Pos)$ is the \emph{output} function.

\paragraph{State encoding: match goals.}
States are encoded as non-empty sets of \emph{match goals}. Write
$\mathit{sub}(\ell) = \{\, \ell|_{p} \mid p \in
\dom_{\setminus\edges}(\ell) \,\}$ for the non-variable subpatterns of
$\ell$, and $\mathit{sub}(\Lhs) = \bigcup_{\ell \in \Lhs}
\mathit{sub}(\ell)$. The set of \emph{match obligations} is
$\mathit{MO} = \mathcal{P}(\mathit{sub}(\Lhs) \times \Pos) \setminus
\{\emptyset\}$, and the set of \emph{match announcements} is
$\mathit{MA} = \Lhs \times \Pos$. A \emph{match goal} is a pair
$(\mathit{mo}, \mathit{ma}) \in \mathit{MO} \times \mathit{MA}$, written
\[
  \ell_{1}@p_{1}, \ldots, \ell_{n}@p_{n} \;\hookrightarrow\; \ell@p,
\]
and read: ``to announce a match for $\ell$ at position $p$, it remains
to observe each subpattern $\ell_{i}$ at position $p_{i}$''. A goal of
the form $\ell@p \hookrightarrow \ell@p$ is \emph{fresh}; one of the form
$\mathit{mo} \hookrightarrow \ell@\epsilon$ is a \emph{root goal}. States
are non-empty subsets of $\mathit{MO} \times \mathit{MA}$, and the
initial state collects the fresh root goals for every left-hand side:
\[
  s_{0} \;=\; \{\, \ell@\epsilon \hookrightarrow \ell@\epsilon \mid \ell \in \Lhs \,\}.
\]
The state-labelling $L$ must select, for each state $s$, a position
appearing in some root goal of $s$; this forces $L(s_{0}) = \epsilon$.
In general the choice is not unique; we fix any canonical policy
(\S\ref{sec:caveats}).

\paragraph{Derivative.}
Let $\mathsf{pos}(\mathit{mo}) = \{\, q \mid \ell@q \in \mathit{mo}
\,\}$. Given a state $s$, a symbol $f$ with $n = \arity{f}$, and the
position $p = L(s)$, the \emph{reduce} operation on a match obligation
is
\[
  \reduce(\mathit{mo}, f, p) \;=\;
  \{\, \ell @q \in \mathit{mo} \mid q \neq p \,\}
  \;\cup\;
  \{\, \ell|_{i} @\, p \cdot i \mid \ell @p \in \mathit{mo},\;
                                  1 \le i \le n,\;
                                  \ell|_{i} \notin V \,\},
\]
applied when $\ell@p \in \mathit{mo}$ for some $\ell$ with
$\hd(\ell) = f$, and yielding $\emptyset$ if no such $\ell$ exists. The
$f$-\emph{derivative} of $s$ is
\[
  \deriv(s, f) \;=\; \mathit{reduced} \,\cup\, \mathit{unchanged} \,\cup\, \mathit{fresh},
\]
where
\begin{align*}
\mathit{reduced} \;&=\; \{\, \reduce(\mathit{mo}, f, L(s)) \hookrightarrow \mathit{ma}
       \mid \mathit{mo} \hookrightarrow \mathit{ma} \in s,\\[-2pt]
       & \qquad\qquad\qquad \exists \ell @L(s) \in \mathit{mo}.\, \hd(\ell) = f,\;
            \reduce(\mathit{mo}, f, L(s)) \neq \emptyset \,\},\\
\mathit{unchanged} \;&=\; \{\, \mathit{mo} \hookrightarrow \mathit{ma} \in s
       \mid L(s) \notin \mathsf{pos}(\mathit{mo}) \,\},\\
\mathit{fresh} \;&=\; \{\, \ell|_{i} @\, L(s) \cdot i
                          \hookrightarrow \ell|_{i} @\, L(s) \cdot i
       \mid \ell \in \Lhs,\; 1 \le i \le n,\; \ell|_{i} \notin V \,\}.
\end{align*}
\emph{Reduced} goals are those whose obligation refers to $L(s)$ and
whose head there matches $f$; \emph{unchanged} goals do not refer to
$L(s)$; \emph{fresh} goals re-seed redex search at the children of
$L(s)$. The output function records completed reductions:
\[
  \eta(s, f) \;=\; \{\, \mathit{ma} \in \mathit{MA} \mid
       f(x_{1}, \ldots, x_{n}) @\, L(s) \hookrightarrow \mathit{ma} \in s,\;
       x_{1}, \ldots, x_{n} \in V \,\}.
\]

\paragraph{Reading a binder symbol.}
When $f$ is a binder with binding positions $\beta_{f}$
(\S\ref{sec:prelim}), the derivative formulae above apply unchanged:
$\reduce$, $\deriv$, and $\eta$ are stated entirely in terms of
position arithmetic and metavariable identification, neither of which
distinguishes binding from non-binding arguments. The only operational
difference is that successor configurations based at $p \cdot i$ for
$i \in \beta_{f}$ inspect symbols of the opened body
$t_{i}(z_{p \cdot i})$, with the fresh free name $z_{p \cdot i}$
acting as a leaf for any subsequent reads (no head symbol, no
derivative spawn). Configuration trees below a binder therefore have
the same shape as configuration trees below an ordinary function
symbol; $\alpha$-equivalence of the fresh names is invariant of the
construction.

\paragraph{Partitioning by dependency; lifting by greatest common prefix.}
The naive choice $\delta(s, f) = \{(\deriv(s, f), \epsilon)\}$ has
infinitely many states, as positions grow without bound. Two finitising
steps cure this: partition $\deriv(s, f)$ into mutually-independent
classes, and shift each class back to a canonical origin.

The \emph{direct dependency} relation $\rdep$ on match goals is
\[
  \bigl(\mathit{mo}_{1} \hookrightarrow \mathit{ma}_{1}\bigr)
  \;\rdep\;
  \bigl(\mathit{mo}_{2} \hookrightarrow \mathit{ma}_{2}\bigr)
  \quad\Longleftrightarrow\quad
  \mathsf{pos}(\mathit{mo}_{1}) \cap \mathsf{pos}(\mathit{mo}_{2}) \neq \emptyset.
\]
$\rdep$ is reflexive and symmetric but not in general transitive; let
$\sim$ be its transitive closure. Partition $\deriv(s, f)$ into
$\sim$-equivalence classes; write $[\deriv(s, f)]_{\sim}$ for the set of
classes. For a class $K$, all match-announcement and obligation
positions in $K$ share a (possibly empty) greatest common prefix
$\gcp(K)$. The \emph{lift} of $K$ rewrites each position
$\gcp(K) \cdot p'$ appearing in $K$ to $p'$:
\[
  \lift(K) \;=\;
  \{\, \{\ell_{i}@\, p'_{i}\}_{i} \hookrightarrow \ell @\, q'
       \mid \{\ell_{i}@\, \gcp(K) \cdot p'_{i}\}_{i}
            \hookrightarrow \ell @\, \gcp(K) \cdot q' \in K \,\}.
\]
The transition function is then
\[
  \delta(s, f) \;=\; \{\, (\lift(K),\, \gcp(K)) \mid K \in [\deriv(s, f)]_{\sim} \,\},
\]
read: ``from configuration $(s, p)$ observing $f$, advance to
$(\lift(K), p \cdot \gcp(K))$ for each class $K$''. Only finitely many
lifted classes can ever arise, so the construction terminates. We write
$\delta^{*}(s, w)$ for the iterated transition along a sequence
$w \in \mathbb{F}^{*}$ of observations.

\paragraph{Outermost-preserving partition.}
The partition by $\sim$ yields the smallest set automaton sound for
arbitrary strategies. For an outermost strategy it is too fine: it
separates goals one would want kept together to discover an outermost
match before its deeper inner matches~\cite[\S 7]{BouwmanErkens2022}.
The \emph{outermost-preserving} relation $\rop$ enlarges $\rdep$ by
also relating goals whose match-announcement positions are
prefix-comparable:
\[
  g_{1} \;\rop\; g_{2}
  \quad\Longleftrightarrow\quad
  g_{1} \;\rdep\; g_{2}
  \;\text{or}\;
  p_{1} \le p_{2}
  \;\text{or}\;
  p_{2} \le p_{1},
\]
where $g_{i} = \mathit{mo}_{i} \hookrightarrow \ell_{i} @\, p_{i}$.
Replacing $\sim$ by the transitive closure of $\rop$ throughout the
construction yields a (typically larger but still finite) automaton in
which a depth-first traversal of the configuration tree yields
outermost matches first. We write $M_{\rdep}$ and $M_{\rop}$ for the
two automata when the distinction matters.

\paragraph{Configurations and configuration trees.}
Running $M$ on a ground term $t$ traverses positions of $t$ under the
guidance of $L$ and $\delta$. A \emph{configuration} is a pair $(s, p)
\in \States \times \Pos$; from it the procedure inspects the symbol
$\hd(t|_{p \cdot L(s)})$ to follow $\delta$. The progress of a
traversal is recorded in a \emph{configuration tree}
\[
  \mathit{CT} \;::=\; \mathsf{B}(s, p) \;\mid\; \mathsf{N}(s, p, \mathit{cts}),
\]
where a \emph{bud} $\mathsf{B}(s, p)$ is an unexplored configuration
and a \emph{node} $\mathsf{N}(s, p, \mathit{cts})$ is an explored
configuration with a finite (possibly empty) set $\mathit{cts}$ of
child configuration trees. The set of nodes of $\mathit{ct}$ is
\[
  \mathsf{nodes}(\mathsf{B}(s, p)) = \emptyset, \qquad
  \mathsf{nodes}(\mathsf{N}(s, p, \mathit{cts}))
  = \{(s, p)\} \cup \!\!\bigcup_{\mathit{ct}' \in \mathit{cts}}\!\!
       \mathsf{nodes}(\mathit{ct}').
\]
The \emph{completed configuration tree} of a ground term $t$ is
\[
  \mathsf{completed}(s, p, t)
  \;=\;
  \mathsf{N}\!\bigl(s,\, p,\, \{\, \mathsf{completed}(s', p \cdot p', t)
                                   \mid (s', p') \in \delta(s, \hd(t|_{p \cdot L(s)})) \,\}\bigr),
\]
with $\mathsf{completed}(t) = \mathsf{completed}(s_{0}, \epsilon, t)$.
The \emph{grow} and \emph{prune} operations are
\begin{align*}
\mathsf{grow}(\mathsf{B}(s, p), t) \;&=\;
  \mathsf{N}\!\bigl(s,\, p,\, \{\, \mathsf{B}(s', p \cdot p')
                                   \mid (s', p') \in \delta(s, \hd(t|_{p \cdot L(s)})) \,\}\bigr),\\
\mathsf{prune}(\mathsf{N}(s, p, \mathit{cts}), q) \;&=\;
  \begin{cases}
    \mathsf{B}(s, p) & \text{if } p \cdot L(s) = q,\\
    \mathsf{N}(s, p, \{\, \mathsf{prune}(\mathit{ct}', q) \mid \mathit{ct}' \in \mathit{cts}\,\}) & \text{otherwise},
  \end{cases}\\
\mathsf{prune}(\mathsf{B}(s, p), q) \;&=\; \mathsf{B}(s, p).
\end{align*}
For a redex at position $q$ inside $t$, there is a unique node of
$\mathsf{completed}(t)$ --- the \emph{initialisation configuration} of
the redex --- whose position is $q$; $\mathsf{prune}(\cdot, q)$
collapses that node back to a bud.

\paragraph{Three results we will lean on.}
The following statements are proved in~\cite[Thm.~1, Thm.~2,
Lemma~1]{BouwmanErkens2022}. We restate them so that
\S\ref{sec:optimality} can cite them locally.

\begin{theorem}[Symbol-once bijection {\cite[Thm.~1]{BouwmanErkens2022}}]\label{thm:be-sym-once}
For any set automaton $M$ for $\Lhs$ and any ground term $t$, the map
\[
  \mathsf{N}(s, p, \mathit{cts}) \;\longmapsto\; p \cdot L(s)
\]
is a bijection between the nodes of $\mathsf{completed}(t)$ and the
domain $\dom(t)$. In particular,
$|\mathsf{nodes}(\mathsf{completed}(t))| = |\dom(t)|$.
\end{theorem}

\begin{theorem}[Matching correctness {\cite[Thm.~2]{BouwmanErkens2022}}]\label{thm:be-match-correct}
Define $\mathsf{matches}(s, p, t) = \{\, (\ell @\, p \cdot q) \mid
(\ell, q) \in \eta(s, \hd(t|_{p \cdot L(s)})) \,\}$. The set of all
redexes of a ground term $t$ for $\Lhs$ is
\[
  \bigcup \{\, \mathsf{matches}(s, p, t) \mid (s, p) \in
              \mathsf{nodes}(\mathsf{completed}(t)) \,\}.
\]
\end{theorem}

\begin{lemma}[Prune properties {\cite[Lemma~1]{BouwmanErkens2022}}]\label{lem:be-prune}
Let $\sqsubseteq$ be the \emph{fragment-of} order on configuration
trees, generated by $\mathsf{B}(s, p) \sqsubseteq \mathsf{N}(s, p,
\mathit{cts})$ and closed componentwise on children. Then for any
$\mathit{ct} \sqsubseteq \mathit{ct}'$ and any position $q$:
\begin{enumerate}[label=\textup{(\roman*)},leftmargin=2.2em]
  \item $\mathsf{prune}(\mathit{ct}, q) \sqsubseteq \mathit{ct}$
        (pruning never adds nodes);
  \item $\mathsf{prune}(\mathit{ct}, q) \sqsubseteq \mathsf{prune}(\mathit{ct}', q)$
        (pruning is monotone in the fragment order);
  \item if $t \xrightarrow{(\ell \rew r) @ q} t'$, then
        $\mathsf{prune}(\mathsf{completed}(t), q)
         = \mathsf{prune}(\mathsf{completed}(t'), q)$
        (firing preserves the trimmed tree); and
  \item the nodes of $\mathit{ct}$ strictly outside the subtree rooted
        at $q$ are preserved by $\mathsf{prune}(\cdot, q)$.
\end{enumerate}
\end{lemma}

\section{The translation scheme}\label{sec:translation}

\subsection{Compilation of terms, rules, and contexts}

We write $\sem{\cdot}$ for the compilation function from terms to processes
(or to names, depending on context); we leave its base cases on function
symbols and variables abstract, since the optimality argument is
parametric in them. The operative clauses are
\eqref{eq:single-rule}--\eqref{eq:context-rule} from the introduction,
which we restate in a form convenient for the analysis.

\begin{definition}[Compilation of rewrite rules]\label{def:compile}
For a single rule $L \rew R$ on translation channel $t\!\ell$,
\[
  \sem{L \rew R}(t\!\ell)
  \;=\;
  \Recv{\sem{L}}{t\!\ell}{\,\Send{t\!\ell}{\sem{R}}\,}.
\]
For a contextual rewrite,
\[
\begin{array}{@{}l@{}}
  \sem{\,L_{1} \rew R_{1}, \ldots, L_{n} \rew R_{n}
        \,\Rightarrow\, K \rew K'\,} \;=\; \\[2pt]
  \quad \mathsf{let}\;tc = \sem{K}\;\mathsf{in} \\
  \quad\quad \mathsf{for}\,((\sem{L_{1}}, \ldots, \sem{L_{n}}) \Leftarrow tc)\,\{ \\
  \quad\quad\quad \Send{tc}{\sem{K'}(\sem{R_{1}}, \ldots, \sem{R_{n}})} \\
  \quad\quad \}.
\end{array}
\]
\end{definition}

The \emph{polymorphic interpretation} of $\Let{tc}{\sem{K}}{\cdots}$ is
this: $\sem{K}$ is a function only of the surface shape of $K$, not of the
hole-fillers. Holes are deliberately absent from the channel formula; they
are precisely the data communicated on $tc$ when the rule fires.

\subsection{What ``optimal'' means here}\label{ssec:optimal-meaning}

Three properties characterise an optimal translation:

\begin{description}[leftmargin=*]
  \item[(O1) Symbol-once.] Each function symbol of $K$, whenever the rule fires,
        is processed by exactly one \textsf{for}-receive in the translation.
  \item[(O2) Prune-preserves.] When an inner contextual rewrite fires
        \emph{below} the position $p$ of the surrounding $K$, the channels
        and pending receives strictly outside $p$ remain valid; no
        re-establishment of the outer channel structure is needed.
  \item[(O3) Coarsest sound.] If two contexts $K$, $K'$ are
        observationally equivalent --- in the sense that no choice of
        hole-filling distinguishes which rule of $\mathcal{R}$ fires --- they
        share a channel; if they are observationally distinct, they receive
        distinct channels.
\end{description}

A naive choice such as $tc(K) = @K$ (reflect the syntactic context
verbatim) trivially satisfies (O3) but fails (O1): two syntactically distinct
contexts that match the same patterns receive different channels and so
duplicate the for-receive work. Choices that collapse too many contexts,
e.g.\ $tc(K) = @\hd(K)$, satisfy (O1) cheaply but fail (O3). The set
automaton sits exactly at the equivalence quotient that satisfies all
three.

\section{The channel-naming construction}\label{sec:construction}

\subsection{From a context to a state}

\begin{definition}[Surface of a context]\label{def:surface}
Let $K$ be an $n$-ary context with hole positions $H_{K} \subseteq \dom(K)$.
The \emph{surface} of $K$ is the partial term obtained from $K$ by
collapsing each hole to a fresh wildcard: $\surface(K)$ has the same
function-symbol positions as $K$, with $\hd(\surface(K)|_{p}) = \hd(K|_{p})$
for $p \in \dom_{\setminus\edges}(K) \setminus H_{K}$, and is undefined on
$H_{K}$.
\end{definition}

\begin{definition}[Automaton trace of a context]\label{def:trace}
Let $M = (\States, s_{0}, L, \delta, \eta)$ be a set automaton for the
left-hand-side set of $\mathcal{R}$. The \emph{automaton trace} of a context
$K$ is the unique configuration tree $T_{M}(K)$ obtained by running the
matching procedure of~\cite[\S 3]{BouwmanErkens2022} on $\surface(K)$,
suspended at every configuration $(s, p)$ whose inspection position
$p \cdot L(s)$ lies in a hole $H_{K}$.
\end{definition}

The trace $T_{M}(K)$ is a finite tree because $\surface(K)$ is finite and
the automaton is deterministic at every transition step.

\begin{definition}[Channel of a context]\label{def:channel}
Define
\[
  tc(K) \;=\; \reflect{\,T_{M}(K)\,},
\]
the canonical reflection of the suspended automaton trace.
\end{definition}

In the unhybridised case where $K$ has a single hole and the automaton is
deterministic enough that $T_{M}(K)$ collapses to a single suspended
configuration $(s, p)$, the formula simplifies to
$tc(K) = \reflect{\delta^{*}(s_{0}, \surface(K))}$, which is the slogan
quoted in the introduction. The general definition in terms of the
suspended configuration tree is the technically correct one, accommodating
$\rdep$-induced parallelism among independent holes and the look-ahead
phenomenon noted in~\cite[\S 1.1]{BouwmanErkens2022}.

\subsection{The compiled rule, in detail}

Substituting Definition~\ref{def:channel} into Definition~\ref{def:compile},
the contextual rule compiles to
\[
\begin{array}{@{}l@{}}
  \mathsf{let}\;tc = \reflect{T_{M}(K)}\;\mathsf{in} \\
  \quad \mathsf{for}\,((\sem{L_{1}}, \ldots, \sem{L_{n}}) \Leftarrow tc)\,\{ \\
  \quad\quad \Send{tc}{\sem{K'}(\sem{R_{1}}, \ldots, \sem{R_{n}})} \\
  \quad \}.
\end{array}
\]
The for-receive on $tc$ blocks until $n$ names arrive, one per hole. Their
arrival corresponds, automaton-side, to taking the suspended transitions
and continuing matching on the hole-fillers. The send on $tc$ in the body
emits the rewritten right-hand side, ready to be received by another
contextual rule whose hole sits where this one's outer pattern fires.

\begin{example}[Boolean simplification]\label{ex:if-not}
Take $\mathcal{R}$ to be the running example
of~\cite[Ex.~8]{BouwmanErkens2022}:
\begin{align*}
  R_{1} &: \mathit{if}(\true, x, y) \rew x &
  R_{2} &: \mathit{if}(\false, x, y) \rew y \\
  R_{3} &: \mathit{not}(\true) \rew \false &
  R_{4} &: \mathit{not}(\false) \rew \true \\
  R_{5} &: \mathit{not}(\mathit{not}(x)) \rew x &&
\end{align*}
The set automaton has four states $s_{0}, s_{1}, s_{2}, s_{3}$ with
$L(s_{0}) = \epsilon$, $L(s_{1}) = 1$, $L(s_{2}) = 1$, and the transitions
of Table~\ref{tab:setautomatonif} (a transcription
of~\cite[Fig.~3]{BouwmanErkens2022}).

\begin{table}[h]
\centering
\renewcommand{\arraystretch}{1.15}
\begin{tabular}{c|c|c|l|l}
state $s$ & $L(s)$ & symbol $f$ & $\delta(s, f)$ & $\eta(s, f)$ \\\hline
$s_{0}$ & $\epsilon$ & $\mathit{if}$    & $\{(s_{1}, \epsilon),\, (s_{0}, 2),\, (s_{0}, 3)\}$ & $\emptyset$ \\
$s_{0}$ & $\epsilon$ & $\mathit{not}$   & $\{(s_{2}, \epsilon),\, (s_{0}, 1)\}$               & $\emptyset$ \\
$s_{0}$ & $\epsilon$ & $\mathit{true}$  & $\emptyset$                                          & $\emptyset$ \\
$s_{0}$ & $\epsilon$ & $\mathit{false}$ & $\emptyset$                                          & $\emptyset$ \\\hline
$s_{1}$ & $1$        & $\mathit{if}$    & $\{(s_{0}, 1),\, (s_{0}, 1.2),\, (s_{0}, 1.3)\}$    & $\emptyset$ \\
$s_{1}$ & $1$        & $\mathit{true}$  & $\emptyset$                                          & $\{R_{1}@\epsilon\}$ \\
$s_{1}$ & $1$        & $\mathit{false}$ & $\emptyset$                                          & $\{R_{2}@\epsilon\}$ \\\hline
$s_{2}$ & $1$        & $\mathit{not}$   & $\{(s_{2}, 1)\}$                                     & $\{R_{5}@\epsilon\}$ \\
$s_{2}$ & $1$        & $\mathit{true}$  & $\emptyset$                                          & $\{R_{3}@\epsilon\}$ \\
$s_{2}$ & $1$        & $\mathit{false}$ & $\emptyset$                                          & $\{R_{4}@\epsilon\}$ \\
$s_{2}$ & $1$        & $\mathit{if}$    & \multicolumn{2}{l}{leads to $s_{3}$; omitted in~\cite[Fig.~3]{BouwmanErkens2022} and unused below} \\
\end{tabular}
\caption{Transitions of the set automaton for $\mathcal{R}$, transcribed
from~\cite[Fig.~3]{BouwmanErkens2022}. A row $(s, f) \mapsto (\delta, \eta)$
records both the partition output $\delta(s, f) \subseteq \States \times \Pos$
(pairs of a target state and a position offset) and the match-announcement
output $\eta(s, f) \subseteq \Lhs \times \Pos$. Symbols on which no
left-hand-side activity is possible are listed with $\delta = \eta = \emptyset$.}
\label{tab:setautomatonif}
\end{table}
\par\medskip\noindent
Consider the context $K = \mathit{if}(\Box, \false, \true)$ in a
contextual rewrite whose inner rule is $R_{5}$ (so the hole is at position
$1$). Running the automaton on $\surface(K)$:
\[
  (s_{0}, \epsilon) \xrightarrow{\;\;\mathit{if}\;\;}
  (s_{1}, \epsilon) \xrightarrow{\;\;\Box\;\;}
  \text{suspend}.
\]
We have $tc(K) = \reflect{(s_{1}, \epsilon)}$, a single-configuration trace.
The compiled rule listens on this channel for an $\mathit{if}$-headed
sub-context filler, performs whatever inner work is needed to bind $x$, and
sends the result back on $tc$.
\par\medskip\noindent
Now consider the context $K' = \mathit{if}(\Box, e, e)$, where the two
$\Box$-positioned arguments are required to coincide (a non-linear hole
pattern). The set-automaton trace is the same up to position $1$, but the
hole-equality constraint is communicated on a side-channel; see
\S\ref{sec:nonlinear}.
\end{example}

\subsection{Algorithmic construction of $tc$}\label{ssec:algorithm}

We summarise the calculation as a procedure (Construction~\ref{constr:tc}) for
clarity. Pre-conditions: a fixed set automaton $M$ for $\mathcal{R}$ has
been constructed off-line.

\begin{construction}[Computing $tc(K)$]\label{constr:tc}\rm
\begin{enumerate}[leftmargin=2em,label=\arabic*.]
  \item Initialise the configuration tree $T \gets \mathsf{B}(s_{0}, \epsilon)$.
  \item While $T$ contains a bud $\mathsf{B}(s, p)$ with
        $p \cdot L(s) \notin H_{K}$:
    \begin{enumerate}[label=2.\arabic*]
      \item Let $f = \hd(\surface(K)|_{p \cdot L(s)})$.
      \item Replace the bud by
            $\mathsf{N}(s, p, \{\mathsf{B}(s', p \cdot p') \mid (s', p') \in \delta(s, f)\})$.
    \end{enumerate}
  \item Return $\reflect{T}$.
\end{enumerate}
\end{construction}

\noindent
The loop terminates because each iteration either grows a bud into a node
with strictly more nodes-of-context-positions, or transitions into the
empty successor set; the number of iterations is bounded by
$|\dom_{\setminus\edges}(\surface(K))|$, in agreement with
Theorem~\ref{thm:be-sym-once}. The reflection in step~3 is a
syntactic operation on a finite tree of automaton states, hence
computable.

\section{Optimality}\label{sec:optimality}

We now state and prove the three optimality theorems. All three are
transports of theorems from~\cite{BouwmanErkens2022} into the channel
setting.

\subsection{Symbol-once inspection}

\begin{theorem}[Symbol-once]\label{thm:symbol-once}
Let $\rho$ be a rho-calculus term obtained by translating a GSLT
$\mathcal{R}$ via Definition~\ref{def:compile} with $tc(\cdot)$ as in
Definition~\ref{def:channel}. For any reduction sequence of $\rho$ that
fires a contextual rewrite at outer-context $K$, each function symbol of
$K$ is consumed by exactly one
\textsf{for}-receive among the compiled rules.
\end{theorem}

\begin{proof}[Proof sketch]
By Construction~\ref{constr:tc}, $tc(K)$ is determined by the suspended
configuration tree $T_{M}(K)$, whose nodes are in bijection with
$\dom_{\setminus\edges}(\surface(K))$ by
Theorem~\ref{thm:be-sym-once}. Each node corresponds to one application
of step~2.2, which, in the operational reading of the channel scheme,
corresponds to one transition through a sub-channel of $tc$ --- i.e.\
one for-receive in the compiled rule. Conversely, no surface symbol of
$K$ is observed twice, by the deterministic-transition guarantee of the
set-automaton construction. Hence the map (surface symbol of $K$)
$\mapsto$ (for-receive that consumes it) is total and injective, and
its image equals the set of for-receives in the compiled rule.
\end{proof}

\subsection{Prune-preserves-work}

In the imperative implementation of~\cite{BouwmanErkens2022}, when a redex
fires at position $p$ the configuration tree is \emph{pruned} to the
initialisation configuration of the redex; matching done strictly above $p$
is preserved. The rho-calculus analogue is that the channel structure
strictly outside the inner-context position remains live across the inner
firing.

\begin{theorem}[Prune-preserves-work]\label{thm:prune-preserves}
Let $\rho$ contain a compiled outer rule $\sem{\Phi_{\mathrm{outer}}}$ with
context $K$ and channel $tc(K)$, and a compiled inner rule
$\sem{\Phi_{\mathrm{inner}}}$ with context $K_{\mathrm{in}}$ situated at
hole position $p$ of $K$. After firing $\Phi_{\mathrm{inner}}$, the
process state restricted to channel $tc(K)$ and to all sub-channels of
$tc(K)$ that correspond to surface symbols of $K$ at positions $q$ with
$q \not\succeq p$ is unchanged.
\end{theorem}

\begin{proof}[Proof sketch]
Translate the prune-correctness lemma (Lemma~\ref{lem:be-prune}) across
Construction~\ref{constr:tc}. The relevant invariant is clause~(i),
$\mathsf{prune}(\mathit{ct}, p) \sqsubseteq \mathit{ct}$, i.e.\ pruning
never adds nodes; combined with clause~(iv), the nodes outside the
subtree rooted at $p$ survive unchanged. In the channel scheme,
``adding nodes'' would correspond to creating new for-receives outside
the prune position. Since no such creation occurs, the for-receives
outside $p$ are exactly those that existed before $\Phi_{\mathrm{inner}}$
fired. The send in the inner rule's body discharges its message on
$tc(K_{\mathrm{in}})$, which is a sub-channel of $tc(K)$; the outer
for-receive then proceeds to match, finding the same surface symbols
of $K$ outside $p$ as before. No re-establishment of the outer channel
structure is required.
\end{proof}

\subsection{Coarsest sound partition}

The third theorem is the central one. It says that no compilation scheme
that compiles different contexts to different channels can collapse more
context pairs than $tc(\cdot)$ does, on pain of unsound firing.

\begin{theorem}[Coarsest sound]\label{thm:coarsest}
Let $\sim_{\mathrm{op}}$ be the equivalence on contexts induced by the
$\rop$-set-automaton: $K \sim_{\mathrm{op}} K'$ iff
$T_{M_{\rop}}(K) = T_{M_{\rop}}(K')$. Then $\sim_{\mathrm{op}}$ is the
coarsest equivalence on contexts such that, for every GSLT $\mathcal{R}$,
$K \sim_{\mathrm{op}} K'$ implies that for any choice of hole-fillers
$t_{1}, \ldots, t_{n}$, the set of rules of $\mathcal{R}$ that fire at the
outermost position of $K[t_{1}, \ldots, t_{n}]$ equals the set of rules
that fire at the outermost position of $K'[t_{1}, \ldots, t_{n}]$.
\end{theorem}

\begin{proof}[Proof sketch]
Soundness: if $K \sim_{\mathrm{op}} K'$, the suspended traces coincide, so
the suspended automaton states coincide, so by the matching-correctness
theorem (Theorem~\ref{thm:be-match-correct}) the set of redexes
announced by the automaton on $K[t_{i}]$ equals the set on $K'[t_{i}]$,
for every choice of $t_{i}$. The outermost subset coincides because
$\rop$ preserves prefix-comparability of match-announcement positions
(\S\ref{ssec:set-automata}).

Coarsest: suppose $\sim'$ is strictly coarser, i.e.\ there exist
$K \sim' K'$ with $T_{M_{\rop}}(K) \neq T_{M_{\rop}}(K')$. Then the two
traces differ at some configuration, so there is some symbol observation
sequence that induces different state-transitions. By the construction of
$\rop$ (\S\ref{ssec:set-automata}), this difference corresponds to at
least one match goal that is reduced in one trace and not in the other.
Choose hole-fillers $t_{i}$ that complete the partial match in one and
not the other (such fillers exist, since the goal-reduction step expects
a specific head symbol). The resulting outermost-match sets at $K[t_{i}]$
and $K'[t_{i}]$ disagree, contradicting soundness of $\sim'$.
\end{proof}

The use of $\rop$ rather than $\rdep$ in Theorem~\ref{thm:coarsest} is
deliberate: $\rdep$ gives a strictly coarser partition that is sound for
\emph{some} strategy but unsound for outermost firing (the very point
made in~\cite[\S 7]{BouwmanErkens2022}, motivating the introduction of
$\rop$). For non-outermost strategies, the analogous coarsest-sound theorem
holds with $\sim_{\mathrm{op}}$ replaced by the equivalence induced by
the strategy-appropriate dependency relation.

\section{Worked examples}\label{sec:examples}

We now present two extended examples that exhibit the construction in
characteristic settings. The first is the lambda calculus with weak
head reduction, the canonical example of a contextual rewrite. The
second is the rho calculus itself compiled into rho --- a reflective
example that illustrates what the optimality theorems deliver when
the object language and the target language coincide.

\subsection{Lambda calculus with head-context reduction}\label{ssec:lambda}

\paragraph{The GSLT.}
Take the function symbols $\mathit{app}$ (arity~2) and $\mathit{lam}$
(arity~1, with binding position~1 in the sense of the binder apparatus
of \S\ref{sec:prelim}); object-language variables \emph{are} the host
calculus's free names introduced when the binder is opened, so no
separate $\mathit{var}$ constructor is needed. The reduction theory
has one principal rule and one contextual rule:
\begin{align*}
  (\beta) \quad
    & \mathit{app}(\mathit{lam}(M), N) \;\rew\; M[N/x], \\
  (\textsc{Head}) \quad
    & M \rew M' \;\;\Rightarrow\;\;
      \mathit{app}(M, N) \rew \mathit{app}(M', N).
\end{align*}
Here $(\textsc{Head})$ encodes weak head reduction: we may reduce in
the function position of an application, but not in the argument
position and not under a $\mathit{lam}$.

The pattern set for the set automaton is
$\Lhs = \{\,\mathit{app}(\mathit{lam}(M), N)\,\}$ --- a single LHS of
depth two. The metavariable $M$ binds to the opened body
$F(z_{1})$ of a matched $\mathit{lam}(F)$; the rewrite RHS
$M[N/x]$ is then the host-level application $F(N)$. These
substitution details are orthogonal to the channel calculation and
do not figure in the automaton.

\paragraph{The set automaton.}
Tracing the construction of \S\ref{ssec:set-automata}:
\begin{itemize}[leftmargin=*]
  \item $s_{0}$: initial state, $L(s_{0}) = \epsilon$, with the single
        root goal
        $\mathit{app}(\mathit{lam}(M), N)@\epsilon \hookrightarrow
         \mathit{app}(\mathit{lam}(M), N)@\epsilon$.
  \item Reading $\mathit{app}$ at $\epsilon$ reduces this goal to
        $\mathit{lam}(M)@1 \hookrightarrow \ell @\epsilon$ (where
        $\ell = \mathit{app}(\mathit{lam}(M), N)$); the variable $N$
        at position $2$ is discarded as a residual. Adding fresh
        root goals at the children of $\epsilon$ and partitioning by
        $\rdep$ yields two equivalence classes: one containing the
        position-$1$ work, one consisting of a fresh root goal at
        position $2$.
  \item After lifting, we obtain
        $\delta(s_{0}, \mathit{app}) = \{(s_{\mathit{app}}, \epsilon),
                                          (s_{0}, 2)\}$.
        The transition $(s_{0}, 2)$ continues redex search in the
        argument position; the transition $(s_{\mathit{app}},
        \epsilon)$ continues the root-level $\beta$-attempt.
  \item $L(s_{\mathit{app}}) = 1$, the unique obligation position of
        $s_{\mathit{app}}$'s root goal.
  \item Reading $\mathit{lam}$ at position $1$ from $s_{\mathit{app}}$
        completes the goal --- $\mathit{lam}$'s body $M$ is a
        variable, so the reduce operation produces $\emptyset$ --- and
        announces the match in $\eta(s_{\mathit{app}},
        \mathit{lam})$.
\end{itemize}
We do not need the rest of the state graph for the channel calculation.
The crucial fact is the position label $L(s_{\mathit{app}}) = 1$.

\paragraph{Computing the channel.}
The contextual rewrite $(\textsc{Head})$ has outer context
$K = \mathit{app}(\Box, N)$, with hole at position $1$ and the variable
$N$ at position~$2$. Since $N$ is a schema variable that ranges over
arbitrary terms, position~$2$ is also effectively a hole from the
matcher's point of view: no LHS depends on its content.\footnote{This
is the precise meaning of the parenthetical in the theorem of
\S\ref{ssec:optimal-meaning}: the surface of $K$ is what the automaton
\emph{actually} consults, not the syntactic rendering of $K$. Schema
variables that no LHS inspects are surface holes.}
Running Construction~\ref{constr:tc} on $K$:
\begin{enumerate}[leftmargin=2em,label=\arabic*.]
  \item Start with $T \gets \mathsf{B}(s_{0}, \epsilon)$.
  \item Bud $\mathsf{B}(s_{0}, \epsilon)$: the inspection position is
        $\epsilon \cdot L(s_{0}) = \epsilon$, not a hole. Read
        $\mathit{app}$. Replace by
        $\mathsf{N}(s_{0}, \epsilon, \{\mathsf{B}(s_{\mathit{app}},
        \epsilon), \mathsf{B}(s_{0}, 2)\})$.
  \item Bud $\mathsf{B}(s_{\mathit{app}}, \epsilon)$: the inspection
        position is $\epsilon \cdot L(s_{\mathit{app}}) = 1$, a hole.
        Suspend.
  \item Bud $\mathsf{B}(s_{0}, 2)$: the inspection position is $2 \cdot
        L(s_{0}) = 2$, a hole. Suspend.
\end{enumerate}
The suspended trace contains a single non-trivial bud at
$(s_{\mathit{app}}, \epsilon)$ awaiting position $1$, plus a
``redex-search'' bud at $(s_{0}, 2)$ awaiting position $2$. The
canonical reflection $tc(K) = \reflect{T_{M}(K)}$ encodes both. The
compiled rule, eliding the boilerplate, is:
\[
\begin{array}{@{}l@{}}
  \mathsf{let}\;tc = \reflect{T_{M}(K)}\;\mathsf{in} \\
  \quad \mathsf{for}\,(\sem{\mathit{app}(\mathit{lam}(M), N)} \Leftarrow tc)\,\{ \\
  \quad\quad \Send{tc}{\sem{\mathit{app}(M[N/x], N)}} \\
  \quad \}.
\end{array}
\]

\paragraph{What optimality delivers.}
The interesting consequences arrive when we consider a chain of head
$\beta$-redexes such as
\[
  t_{0} \;=\; \mathit{app}(\mathit{app}(\mathit{app}(\mathit{lam}(M_{3}),
              N_{3}), N_{2}), N_{1}),
\]
where weak head reduction proceeds by three successive $\beta$-firings
at progressively deeper positions in the head spine.

\begin{itemize}[leftmargin=*]
  \item \emph{Symbol-once} (Theorem~\ref{thm:symbol-once}). The three
        $\mathit{app}$ symbols of $t_{0}$ are processed by exactly
        three for-receives in the channel encoding, regardless of how
        many $\beta$-firings occur. Naive translation, in which each
        head-context instance allocates fresh receivers for the
        application symbols already on the spine, would re-process the
        outer $\mathit{app}$s after every inner reduction --- giving
        quadratic work in the spine length. The set-automaton encoding
        keeps it linear.
  \item \emph{Prune-preserves-work} (Theorem~\ref{thm:prune-preserves}).
        After the innermost $\beta$ fires --- replacing
        $\mathit{app}(\mathit{lam}(M_{3}), N_{3})$ with $M_{3}[N_{3}/x]$
        --- the for-receives associated with the outer two
        $\mathit{app}$s remain live on their channels. They have already
        consumed their respective $\mathit{app}$ symbols and are now
        blocked waiting for hole content; no re-establishment is
        needed.
  \item \emph{Coarsest-sound} (Theorem~\ref{thm:coarsest}). The argument
        $N$ in $K = \mathit{app}(\Box, N)$ does not appear in any LHS
        of $\Lhs$, so the suspended trace $T_{M}(K)$ is the same for
        every choice of $N$. All instances of $(\textsc{Head})$ share
        a single channel structure regardless of argument shape; this
        is the maximal sharing consistent with the firing semantics.
\end{itemize}

In short, weak head normalisation of a head-redex chain of length $n$
takes $O(n)$ work in the channel encoding, with sharing across instances
of the head-context rule that no syntax-directed translation could
achieve.

\subsection{The rho calculus into itself}\label{ssec:rho-into-rho}

\paragraph{The GSLT.}
We now take the source GSLT to be the rho calculus presented as a term
rewriting theory, and compile it into the rho calculus itself. The
function symbols are the syntactic constructors of rho:
$\mathit{par}$ (arity~2), $\mathit{in}$ (arity~3), $\mathit{out}$
(arity~2), $\mathit{drop}$ (arity~1), $\mathit{nil}$ (arity~0). We
write $\mathit{in}(x, y, P)$ for $\Recv{y}{x}{P}$ and $\mathit{out}(x,
P)$ for $\Send{x}{P}$. The reduction theory has the principal
$\textsc{Comm}$ rule and two structural contextual rules:
\begin{align*}
  (\textsc{Comm}) \quad
    & \mathit{par}(\mathit{in}(x, y, P), \mathit{out}(x, Q))
      \;\rew\; P[@Q / y], \\
  (\textsc{Par}_{L}) \quad
    & M \rew M' \;\;\Rightarrow\;\;
      \mathit{par}(M, Q) \rew \mathit{par}(M', Q), \\
  (\textsc{Par}_{R}) \quad
    & M \rew M' \;\;\Rightarrow\;\;
      \mathit{par}(P, M) \rew \mathit{par}(P, M').
\end{align*}
The $\textsc{Comm}$ rule has a non-linear LHS: the channel name $x$
appears in both the $\mathit{in}$- and the $\mathit{out}$-prefix. We
treat the non-linearity per \S\ref{sec:nonlinear}, with a consistency
receive that checks $x =_{\Names} x'$.

\paragraph{The set automaton.}
The pattern set is $\Lhs = \{\,\mathit{par}(\mathit{in}(x, y, P),
\mathit{out}(x', Q))\,\}$ (linearised by renaming the second occurrence
of $x$, with the equality recorded for non-linear post-checking). We
trace the relevant part of the construction:
\begin{itemize}[leftmargin=*]
  \item $s_{0}$: initial, $L(s_{0}) = \epsilon$, with the fresh root
        goal at $\epsilon$.
  \item Reading $\mathit{par}$ at $\epsilon$ produces the reduced
        obligations $\mathit{in}(x, y, P)@1 \hookrightarrow \ell
        @\epsilon$ and $\mathit{out}(x', Q)@2 \hookrightarrow
        \ell@\epsilon$, plus fresh root goals at positions $1$ and $2$.
  \item Crucially, partitioning by $\rdep$ yields \emph{two}
        equivalence classes: positions $1$ and $2$ are independent.
        Lifting gives a transition with two target states ---
        $(s_{\mathit{in}}, 1)$ for the in-prefix work and
        $(s_{\mathit{out}}, 2)$ for the out-prefix work --- which the
        matcher may explore in either order or in parallel.
  \item Each of $s_{\mathit{in}}$ and $s_{\mathit{out}}$ has its own
        local labelling and transitions; details are unilluminating for
        the channel calculation. The match for the $\textsc{Comm}$
        LHS is announced once both branches reach their respective
        completion states.
\end{itemize}
The non-linearity check $x =_{\Names} x'$ is dispatched, per
Definition~\ref{def:consistency}, on a side channel.

\paragraph{Computing the channel.}
For the contextual rule $(\textsc{Par}_{L})$, the outer context is
$K_{L} = \mathit{par}(\Box, Q)$ with hole at position~$1$ and schema
variable $Q$ at position~$2$. As in the lambda example, $Q$ is a
schema variable that no LHS inspects \emph{at the outer level} ---
the only LHS is the $\textsc{Comm}$ pattern, and that pattern's match
work at position $2$ is initiated only after the entire LHS commits to
matching at the root. From the standpoint of $(\textsc{Par}_{L})$'s
own surface, position~$2$ is a hole. Construction~\ref{constr:tc}
yields:
\begin{enumerate}[leftmargin=2em,label=\arabic*.]
  \item Read $\mathit{par}$ at $\epsilon$, producing buds at
        $(s_{\mathit{in}}, 1)$ and $(s_{\mathit{out}}, 2)$.
  \item Bud $(s_{\mathit{in}}, 1)$: inspection position is $1 \cdot
        L(s_{\mathit{in}})$; if the relevant offset is in the
        hole-region of $K_{L}$ (which it is, the hole being at
        position~$1$), suspend.
  \item Bud $(s_{\mathit{out}}, 2)$: inspection position is $2 \cdot
        L(s_{\mathit{out}})$; this is in the schema-variable region of
        $K_{L}$ at the outer level, suspend.
\end{enumerate}
Both branches are suspended, the two suspensions are independent (they
arose from distinct $\rdep$-classes), and $tc(K_{L}) =
\reflect{T_{M}(K_{L})}$ records both suspended configurations. The
symmetric calculation gives $tc(K_{R})$ for $(\textsc{Par}_{R})$, with
the in/out roles of positions $1$ and $2$ exchanged.

\paragraph{What optimality delivers.}
This example highlights five distinct benefits, three from the
optimality theorems of \S\ref{sec:optimality} and two that are
specific to the reflective setting.

\begin{itemize}[leftmargin=*]
  \item \emph{Symbol-once.} In a process pool $P_{1} \mid P_{2} \mid
        \cdots \mid P_{n}$, parsed as a left-associated $\mathit{par}$
        tree of depth $n-1$, each $\mathit{par}$ is consumed by exactly
        one for-receive. The number of $\mathit{par}$-symbol
        observations in the meta-evaluator's run is $O(n)$, not the
        $O(n^{2})$ of a naive ``try every pair of children'' matcher.
  \item \emph{Concurrent firing.} The $\rdep$-independence of
        positions~$1$ and~$2$ in the $\textsc{Comm}$ LHS means
        $tc(K_{L})$ and $tc(K_{R})$ refer to genuinely independent
        sub-channels. The compiled rules for $(\textsc{Par}_{L})$ and
        $(\textsc{Par}_{R})$ may fire concurrently --- as the rho
        calculus's own structural-congruence rules already permit ---
        and the channel encoding inherits this concurrency directly
        from the automaton's partition structure.\footnote{Recall the
        forward-looking remark of \S\ref{sec:caveats} on parallel
        firing: this example is the cleanest setting in which to
        observe that $\rdep$ is doing the work of a concurrency
        analysis.}
  \item \emph{Prune-preserves-work.} In a deeply nested process tree
        with multiple simultaneously-enabled $\textsc{Comm}$ redexes,
        the firing of any one redex does not invalidate the for-receives
        associated with the surrounding $\mathit{par}$ structure. The
        meta-evaluator does not have to ``re-parse'' the process tree
        after each communication step.
  \item \emph{Coarsest-sound.} All contexts $\mathit{par}(\Box, Q)$
        share the channel $tc(K_{L})$ regardless of $Q$, and dually for
        $(\textsc{Par}_{R})$. The redex-discovery work is shared across
        every site where the corresponding contextual rule could fire.
  \item \emph{Reflective coherence.} The meta-evaluator's channel
        structure mirrors the object-level rho calculus's channel
        structure: outer $\mathit{par}$s correspond to outer channels,
        inner $\textsc{Comm}$ redexes correspond to inner channels, and
        substitution is realised by sending and receiving names. The
        optimality of the matcher is therefore inherited as the
        optimality of the meta-evaluator's parallel structure ---
        precisely the coherence property that makes a self-applicable
        rewrite-based language definition (in the sense
        of~\cite{F1R3FLY-mettail}) admit efficient bootstrapping.
\end{itemize}

A pleasant corollary: in a process pool with $n$ parallel components and
$m$ enabled $\textsc{Comm}$ redexes, the meta-evaluator's redex-discovery
phase costs $O(n + m)$ communications, not the $O(n^{2})$ that a naive
self-interpreter would incur from trying every $(\mathit{in},
\mathit{out})$ pair. The optimality of the underlying set automaton is
exactly what bridges the gap.

\section{Non-linear extensions}\label{sec:nonlinear}

When some $L_{i}$ is non-linear --- some variable repeats --- the set
automaton is not powerful enough to enforce the equality constraint:
no regular tree automaton can recognise the language $\{t \mid t|_{p}
\equiv t|_{q}\}$ for two fixed positions $p \neq q$~\cite{Comon2007}.

\paragraph{Non-linear partition and consistency.}
For a pattern $\ell$, the \emph{non-linear partition}
$\pi(\ell) \subseteq \mathcal{P}(\edges(\ell))$ is the partition of the
variable positions of $\ell$ by variable identity: $p, q \in \edges(\ell)$
lie in the same class of $\pi(\ell)$ iff $\ell|_{p}$ and $\ell|_{q}$
are the same variable. A pattern is linear iff every class of
$\pi(\ell)$ is a singleton; non-linearity is the presence of a class of
size $\ge 2$. A term $t$ is \emph{consistent} with $\pi(\ell)$ at
position $p \in \dom(t)$ iff for every class $P \in \pi(\ell)$ and
every $q_{1}, q_{2} \in P$, $t|_{p \cdot q_{1}} \equiv t|_{p \cdot
q_{2}}$. A linear pattern is vacuously consistent.

\paragraph{Pre-matches and the ambiguous/enabled/disabled trichotomy.}
A \emph{pre-match} is a pair $(\ell, p)$ such that $t|_{p}$ matches
$\ell$ up to variable identification (i.e.\ $t|_{p}$ matches the
linearisation of $\ell$). Bouwman--Erkens delegate consistency
checking to the rewrite procedure, classifying every pre-match as
\emph{ambiguous} (consistency unchecked), \emph{enabled} (consistency
verified), or \emph{disabled} (consistency violated); only enabled
pre-matches may fire, and inner firings move pre-matches whose
dependent subterms have changed back to \emph{ambiguous}, where they
will be re-checked~\cite[Algorithm~2]{BouwmanErkens2022}. We mirror
this in the channel scheme by introducing auxiliary \emph{consistency
channels}.

\begin{definition}[Consistency receive]\label{def:consistency}
Let $L$ be a pattern with non-linear partition $\pi(L)$. For each
non-singleton class $P \in \pi(L)$ with $P = \{q_{1}, \ldots, q_{k}\}$
($k \ge 2$), define a guarded receive
\[
\begin{array}{@{}l@{}}
  \Match{P}(t) \;=\; \\
  \quad \mathsf{for}\,(z_{1}, \ldots, z_{k} \Leftarrow @t|_{q_{1}}, \ldots, @t|_{q_{k}})\,\{ \\
  \quad\quad \mathit{if}\;\bigwedge_{i=2}^{k} z_{1} =_{\mathcal{N}} z_{i}\;\mathit{then}\;\mathit{enable}\;\mathit{else}\;\mathit{disable} \\
  \quad \}
\end{array}
\]
where $=_{\mathcal{N}}$ is the rho-calculus name equality of
\S\ref{sec:prelim}.
\end{definition}

The compiled non-linear rule then becomes
\[
\begin{array}{@{}l@{}}
  \mathsf{let}\;tc = \reflect{T_{M}(K)}\;\mathsf{in} \\
  \quad \mathsf{for}\,((\sem{L_{1}}, \ldots, \sem{L_{n}}) \Leftarrow tc)\,\bigl\{ \\
  \quad\quad \bigl(\bigparallel_{i, P \in \pi(L_{i}), |P| \ge 2} \Match{P}(L_{i})\bigr) \\
  \quad\quad \mid \mathit{enable\text{-}guarded}(\Send{tc}{\sem{K'}(\ldots)}) \\
  \quad \bigr\},
\end{array}
\]
where $\mathit{enable\text{-}guarded}(P)$ denotes the process that
proceeds as $P$ once one $\mathit{enable}$ token has been received per
non-singleton class on a dedicated coordination channel, and remains
blocked otherwise. The send to fire the rule is emitted only if all
consistency receives report \emph{enable}. This is exactly the
\textit{ambiguous} $\to$ \textit{enabled} transition
of~\cite[Algorithm~2, lines 12--17]{BouwmanErkens2022}, lifted to
channel form. The \textit{disabled} state holds via the absence of the
enable signal.

\begin{remark}[Re-checking on inner firing]\label{rem:recheck}
\cite[Algorithm~2, lines 24--31]{BouwmanErkens2022} re-classifies a
disabled or enabled match as ambiguous if an inner rewrite changes a
subterm referenced by a non-singleton class of $\pi(L)$. In the channel
scheme, the analogous behaviour is automatic: on inner firing, the inner
rule's send pushes a fresh value to $tc(K_{\mathrm{in}})$, which causes
the enclosing $\Match{P}$ receives to re-fire with the new value. No
explicit administrative loop is required, modulo the persistent-receive
semantics of $\Leftarrow$.
\end{remark}

\section{Caveats and future work}\label{sec:caveats}

\paragraph{Look-ahead.}
The set automaton requires a look-ahead bounded by the maximum LHS depth
\cite[\S 1.1]{BouwmanErkens2022}; equivalently, the channel $tc(K)$ may
suspend before observing the entire surface of $K$, waiting for hole
content to determine which sub-pattern to match. This corresponds, in
implementation terms, to ``intermediate channels'' that exist between
$tc(K)$ and $tc(K_{\mathrm{in}})$. Their construction is mechanical
(unfold step~2 of Construction~\ref{constr:tc} until either a match or a
hole is reached) but they should be made explicit in any concrete
implementation, as they bear on memory cost.

\paragraph{Canonical state choice.}
The state-labelling function $L : \States \to \Pos$ is, in
\cite{BouwmanErkens2022}, chosen non-deterministically among the
positions of root match obligations. For $tc(\cdot)$ to be a function
rather than a relation, the choice must be fixed at automaton-construction
time; a lexicographic policy on positions is the obvious canonical
choice and is what one would implement.

\paragraph{Dynamic rule sets.}
A fully-reflective GSLT --- e.g.\ MeTTaIL~\cite{F1R3FLY-mettail} ---
permits new rewrite rules to be introduced at runtime. The
construction of \S\ref{sec:construction} assumes a closed set $\Lhs$ to
build the automaton from. Two extension strategies are available:
incremental automaton update (recompute $\delta$ for the affected states
as new patterns are added) or stratified compilation (a fresh
sub-automaton per dynamic rule batch, dispatched by an outer name-equality
check). The former preserves the optimality theorems but at higher
construction cost; the latter is cheaper but degrades (O3) to a coarser
partition that may distinguish a context $K$ across batches.

\paragraph{Concurrent firing.}
The rho calculus permits concurrent application of non-conflicting rules
in a way that the imperative model of~\cite{BouwmanErkens2022} does not.
The $\rdep$-partition at the heart of the set-automaton construction is
exactly an analysis of independent sub-matchings; under the channel
encoding it becomes an analysis of independent sub-channels that may be
serviced in parallel. A precise statement and proof of a
parallel-firing theorem --- ``two contextual rewrites whose outer contexts
share no $\rdep$-class fire concurrently'' --- is a natural follow-up to
this paper.

\paragraph{Right-hand-side precompilation.}
\cite[\S 8]{BouwmanErkens2022} suggest precompiling the configuration
tree fragments arising from right-hand sides, so that the symbols of $R$
need not be matched again after substitution. The channel-side analogue
is to precompile the trace $T_{M}(R)$ (with holes at the variables of
$R$) and have the inner rule's send emit not $\sem{R}$ but a process
already partially matched, attached at the trace's bud channels. The
optimisation is an adaptation of Construction~\ref{constr:tc} to
right-hand sides and is straightforward.

\section{Conclusion}\label{sec:conclusion}

We have given a closed-form construction of the channel $tc(K)$ in the
contextual-rewrite translation scheme of Definition~\ref{def:compile},
identified it with the canonical reflection of a set-automaton trace, and
proved that it is optimal in three precise senses inherited from the
Bouwman--Erkens framework. The construction is parametric in the state
encoding, in the choice of dependency relation (and hence of strategy),
and in the rho-calculus reflection bracket; this parametricity is what
makes the polymorphic $\Let{tc}{\sem{K}}{\cdots}$ step --- the original
locus of the question --- precise, computable, and theory-grounded. The
non-linear extension and the dynamic-rule extension are sketched and left
as engineering follow-ups.

\paragraph{Acknowledgements.}
The author thanks Mike Stay, Brian Beckman, and Hannah Adams for ongoing
discussions that shaped the F1R3FLY platform context in which this work
sits, and acknowledges the Bouwman--Erkens paper~\cite{BouwmanErkens2022}
as the immediate technical foundation.

\begin{thebibliography}{9}

\bibitem{BouwmanErkens2022}
M.~Bouwman and R.~Erkens.
\newblock Term rewriting based on set automaton matching.
\newblock arXiv:2202.08687, 2022.

\bibitem{ErkensGroote2021}
R.~Erkens and J.F.~Groote.
\newblock A set automaton to locate all pattern matches in a term.
\newblock In \emph{Theoretical Aspects of Computing -- ICTAC 2021},
  LNCS 12819, Springer, pp.~67--85, 2021.

\bibitem{Baader1998}
F.~Baader and T.~Nipkow.
\newblock \emph{Term Rewriting and All That}.
\newblock Cambridge University Press, 1998.

\bibitem{MeredithRadestock2005}
L.G.~Meredith and M.~Radestock.
\newblock A reflective higher-order calculus.
\newblock \emph{Electronic Notes in Theoretical Computer Science},
  141(5):49--67, 2005.

\bibitem{Comon2007}
H.~Comon, M.~Dauchet, R.~Gilleron, C.~L\"oding, F.~Jacquemard,
  D.~Lugiez, S.~Tison, and M.~Tommasi.
\newblock Tree automata techniques and applications.
\newblock \texttt{http://tata.gforge.inria.fr/}, 2007.

\bibitem{F1R3FLY-mettail}
L.G.~Meredith.
\newblock The MeTTaIL language: defining languages as triples of syntactic
  types, equations, and rewrites.
\newblock F1R3FLY.io technical report and prototype,
  \texttt{https://github.com/F1R3FLY-io/mettail-rust}, 2025.

\end{thebibliography}

\end{document}
